%\documentclass{article}
\documentclass[notitlepage,reqno,11pt]{amsart}


\usepackage{graphicx} % Required for inserting images
\usepackage{amsmath}
\usepackage{amsfonts}
\usepackage{mathtools}
\usepackage{bbm}

% \newcommand{\mb}[1]{\mathbb{#1}}
\newcommand{\F}{\mathcal{F}}
\newcommand{\ld}{\lambda}
\newcommand{\s}{\sigma}
\newcommand{\Pro}{\mathbb{P}}
\newcommand{\E}{\mathbb{E}}
\newcommand{\id}{\mathbbm{1}}
\newcommand{\de}{\text{d}}



%\usepackage[hmarginratio=1:1]{geometry}
%\usepackage{amsthm}
%\usepackage{showlabels}
%\usepackage{showkeys}
\usepackage{latexsym,amssymb, epsfig, amsmath,amsfonts, subfigure,amsthm,mathrsfs}
\renewcommand{\labelenumi}{(\roman{enumi})}
\usepackage{rotating}
\usepackage[toc,page]{appendix}
\usepackage{color}
\usepackage{multirow}
\usepackage{relsize}
\usepackage{microtype}
\usepackage[foot]{amsaddr}
\usepackage{dsfont}

%\usepackage{ulem}
%\setlength{\textwidth}{6.5in}
%\setlength{\textheight}{9in}
%\setlength{\oddsidemargin}{.0in}
%\setlength{\topmargin}{-0.3in}
%\setlength{\headheight}{0in}

%\usepackage[colorlinks,anchorcolor=blue,citecolor=blue]{hyperref}

\usepackage[colorlinks, allcolors=blue]{hyperref}

%\usepackage{ulem}

\usepackage[margin=1in]{geometry}



%\topmargin 0.0truein
%\oddsidemargin 0.25truein
%\evensidemargin 0.25truein
%\textheight 8.5truein
%\textwidth 6.0truein

%\raggedbottom


\numberwithin{equation}{section}
 \newtheorem{assumption}{Assumption}[section]
\newtheorem{lemma}{Lemma}[section]
\newtheorem{theorem}{Theorem}[section]
\newtheorem{definition}{Definition}[section]
\newtheorem{coro}{Corollary}[section]
\newtheorem{prop}{Proposition}[section]
\newtheorem{remark}{Remark}[section]
\newtheorem{condition}{Condition}[section]
\newtheorem{example}{Example}[section]
\newtheorem{conjecture}{Conjecture}[section]

\newlength{\defbaselineskip}
\setlength{\defbaselineskip}{\baselineskip}
\newcommand{\setlinespacing}[1]%
           {\setlength{\baselineskip}{#1 \defbaselineskip}}
\newcommand{\doublespacing}{\setlength{\baselineskip}%
                           {1.5 \defbaselineskip}}
\newcommand{\singlespacing}{\setlength{\baselineskip}{\defbaselineskip}}




\newcommand{\bsq}{\vrule height .9ex width .8ex debth -.1ex}
\newcommand{\CC}{{\mathbb C}}
\newcommand{\DD}{{\mathbb D}}
\newcommand{\RR}{{\mathbb R}}
\newcommand{\ZZ}{{\mathbb Z}}
\newcommand{\NN}{{\mathbb N}}
\newcommand{\ME}{{\mathcal{E}^{(N)}_0}}
\newcommand{\MC}{{\mathcal{C}^{(N)}_0}}
\newcommand{\MD}{{\mathcal{D}^{(N)}_0}}
\def\E{\mathbb{E}}
\def\P{\mathbb{P}}
\def\R{\mathbb{R}}
\def\wconv{\stackrel{w}{\rightarrow}}
\newcommand{\sE}{{\mathcal{E}}}
\newcommand{\sT}{{\mathcal{T}}}
\newcommand{\sA}{{\mathcal{A}}}
\newcommand{\sD}{{\mathcal{D}}}
\newcommand{\sB}{{\mathcal{B}}}
\newcommand{\sC}{{\mathcal{C}}}
\newcommand{\sG}{{\mathcal{G}}}
\newcommand{\sF}{{\mathcal{F}}}
\newcommand{\sI}{{\mathcal{I}}}
\newcommand{\sJ}{{\mathcal{J}}}
\newcommand{\sH}{{\mathcal{H}}}
\newcommand{\sP}{{\mathcal{P}}}
\newcommand{\sN}{{\mathcal{N}}}
\newcommand{\sR}{{\mathcal{R}}}
\newcommand{\sS}{{\mathcal{S}}}
\newcommand{\sL}{{\mathcal{L}}}
\newcommand{\sM}{{\mathcal{M}}}
\newcommand{\sX}{{\mathcal{X}}}
\newcommand{\sZ}{{\mathcal{Z}}}
\newcommand{\hsp}{\hspace*{\parindent}}
\newcommand{\lf}{\lfloor}
\newcommand{\rf}{\rfloor}
\newcommand{\lc}{\lceil}
\newcommand{\rc}{\rceil}
\newcommand{\ra}{\rightarrow}
\newcommand{\deq}{\stackrel{\rm d}{=}}
\newcommand{\beql}[1]{\begin{equation}\label{#1}}
\newcommand{\eeq}{\end{equation}}
\newcommand{\eqn}[1]{(\ref{#1})}
\newcommand{\beqal}[1]{\begin{eqnarray}\label{#1}}
\newcommand{\eeqa}{\end{eqnarray}}
\newcommand{\beq}{\begin{displaymath}}
\newcommand{\eeqno}{\end{displaymath}}
\newcommand{\bali}[1]{\begin{align}\label{#1}}
\newcommand{\eali}{\begin{align}}
\newcommand{\balino}{\begin{align*}}
\newcommand{\ealino}{\begin{align*}}
\newcommand{\vsp}{\vspace{.1in}}
\newcommand{\dd}{\ldots}
\newcommand{\af}{\alpha}
\newcommand{\bt}{\beta}
\newcommand{\dt}{\delta}
\newcommand{\la}{\lambda}
\newcommand{\tht}{\theta}
\newcommand{\gm}{\gamma}
\newcommand{\kp}{\kappa}
\newcommand{\ep}{\varepsilon}
\newcommand{\La}{\Lambda}
\newcommand{\bA}{{\mathbf A}}

\newcommand{\rB}{{\mathrm B}}
\newcommand{\rR}{{\mathrm R}}

\newcommand{\Var}{\text{\rm Var}}
\newcommand{\Cov}{\text{\rm Cov}}
\newcommand{\wt}{\widetilde}

\newcommand{\D}{{\mathrm d}}

\newcommand{\mfI}{\mathfrak{I}}
\newcommand{\mfi}{\mathfrak{i}}
\newcommand{\mfE}{\mathfrak{E}}
\newcommand{\mfF}{\mathfrak{F}}
\newcommand{\mfe}{\mathfrak{e}}
\newcommand{\mfR}{\mathfrak{R}}
\newcommand{\mfr}{\mathfrak{r}}
\newcommand{\mfa}{\mathfrak{a}}
\newcommand{\eps} {\varepsilon}
\newcommand{\bzero}{{\mathbf 0}}
\newcommand{\bX}{{\mathbf X}}
\newcommand{\bY}{{\mathbf Y}}
\newcommand{\bD}{{\mathbf D}}
\newcommand{\bL}{{\mathbf L}}
\newcommand{\bS}{{\mathbf S}}
\newcommand{\bs}{{\mathbf s}}
\newcommand{\bN}{{\mathbf N}}
\newcommand{\bn}{{\mathbf n}}
\newcommand{\bC}{{\mathbf C}}
\newcommand{\bu}{{\mathbf u}}
\newcommand{\bc}{{\mathbf c}}
\newcommand{\bQ}{{\mathbf Q}}
\newcommand{\bB}{{\mathbf B}}
\newcommand{\bF}{{\mathbf F}}
\newcommand{\bG}{{\mathbf G}}
\newcommand{\bH}{{\mathbf H}}
\newcommand{\bR}{{\mathbf R}}
\newcommand{\bV}{{\mathbf V}}
\newcommand{\bW}{{\mathbf W}}
\newcommand{\bI}{{\mathbf I}}
\newcommand{\bP}{{\mathbf P}}
\newcommand{\bfe}{{\bf e}}
\newcommand{\bK}{{\mathbf K}}
\newcommand{\bM}{{\mathbf M}}
\newcommand{\bU}{{\mathbf U}}
\newcommand{\bT}{{\mathbf T}}
\newcommand{\ua}{\uparrow}
\newcommand{\bZ}{{\mathbf Z}}
\newcommand{\bz}{{\mathbf z}}
\newcommand{\by}{{\mathbf y}}
\newcommand{\bw}{{\mathbf w}}
\newcommand{\bq}{{\mathbf q}}
\newcommand{\bb}{{\mathbf b}}
\newcommand{\ba}{{\mathbf a}}
\newcommand{\bm}{{\mathbf m}}
\newcommand{\bx}{{\mathbf x}}
\newcommand{\be}{{\mathbf e}}
\newcommand{\bone}{{\mathbf 1}}
\newcommand{\hv}{\widehat{v}}
\newcommand{\hV}{\widehat{V}}
\newcommand{\hf}{\widehat{f}}
\newcommand{\hg}{\widehat{g}}
\newcommand{\qandq}{\quad\mbox{and}\quad}
\newcommand{\qifq}{\quad\mbox{if}\quad}
\newcommand{\qiffq}{\quad\mbox{if and only if}\quad}
\newcommand{\qforq}{\quad\mbox{for}\quad}
\newcommand{\qorq}{\quad\mbox{or}\quad}
\newcommand{\qasq}{\quad\mbox{as}\quad}
\newcommand{\qinq}{\quad\mbox{in}\quad}
\newcommand{\qforallq}{\quad\mbox{for all}\quad}
\newcommand{\df}{\displaystyle\frac}
\newcommand{\non}{\nonumber}
\newcommand{\RA}{\Rightarrow}
\newcommand{\ti}{\tilde}
\newcommand{\ch}{\check}
\newcommand{\bref}[1]{(\ref{#1})}
\newcommand{\baa}{\begin{eqnarray*}}
\newcommand{\eaa}{\end{eqnarray*}}
\def\binom#1#2{{#1}\choose{#2}}

\newcommand{\mds}{\mathds{1}}

\begin{document}

\newcommand{\ttl}{\Large Proofs for Boros}

%\newcommand{\ttls}{ }


\title[]{\ttl}



\maketitle

\section{Notation}
% Denote the set of users ${U^{(i)}}_{i\in[1,...,N]}$. Each user’s state is a tuple 
% $$
% U^{(i)}(t) = (c^{(i)}(t),s^{(i)}(t),s^{(i)}_k(t))_{k\in[-K,...,K]}
% $$ where $c$ is the collateral, $s_k$ is the position size, the size of the limit order price on tick $k$, where $s_k=0$ indicates no order in that tick rate. 

We let $r_k$ denotes the corresponding annualized rate at tick $k$ such that $r_k = r(k)\times 365\times 3$ where $r(k)$ is the real fractional payment of each payment, e.g. 0.0001 in funding rate.

We use time to maturity $t$ with the unit in years, e.g. $t=1/4$ correspond with 3 months until maturity. Note that this $t$ is discretized, meaning it takes value $\{0, \frac{1}{C}, \frac{2}{C}, ...\}$.

The environment includes the highest bid tick $bid(t)$, the lowest ask price $ask(t)$, and the mark price $m(t)$ which is “close” to the mid price in some sense. 
\begin{align*}
    bid(t) &= \max_k\{k:s^{(i)}_k(t)>0 \text{ for some } i\}\\
    ask(t) &= \min_k\{k:s^{(i)}_k(t)<0 \text{ for some } i\}\\
\end{align*}

% \section{State change}

% Each user’s state only changes during these events:
% \begin{itemize}
%     \item Realize PnL between time windows: $c\to c+\dt c$
%     \item Open limit order at tick $k$: $s_k\to s_k + \dt s$ where $s_k\times \dt s > 0$. Requirement: if $\dt s > 0$ then $k<ask$, otherwise $k > bid$.
%     \item Cancel limit order at tick $k$: $s_k\to s_k - \dt s$ where $s_k\times \dt s > 0$. Requirement: $s_k \neq 0$ and $s_k(s_k - \dt) \geq 0 $
%     \item Limit Order filled at tick $k$: 
%     \begin{align*}
%         s_k \to s_k-\dt s\\
%         s\to s + \dt s\\
%         c \to c - \dt c
%     \end{align*}
%     where $s_k\times \dt s > 0$, $\dt c = \dt s\times t\times r_k$. 
%     \item Market order at tick $k$: 
%     \begin{align*}
%         s\to s - \dt s\\
%         c \to c + \dt c
%     \end{align*}
%     where $\dt c = \dt s\times t\times r_k$.
    
%     Requirement: if $\dt s > 0$ then $k=bid$, otherwise $k=ask$.
% \end{itemize}

% Limit Order filled and Market order happen between 2 users only as a pair of simultaneous actions at a specific tick rate $k$ and time $t$, i.e. there exists a pair of users $i,j$ such that:
% \begin{align*}
%     s^{(i)}_k \to s^{(i)}_k-\dt s,s^{(i)}\to s^{(i)} + \dt s\\
%     s^{(j)}\to s^{(j)}-\dt s
% \end{align*}

\section{Margin Requirement}

We revisit the Margin system here for reader's convenience. This is the version for one market.

\subsection{Initial Margin}

Each limit order is a tuple $(k, s_k)$ where $k$ is the tick and $s_k$ is the order size. Note that at time $t$, $s_k> 0\RA k\leq bid$ and $s_k< 0\RA k\geq ask$. Such an open order has the (pre-scale) initial margin 
$$PIM_k = |s_k|\times\max\{I,|r_k|\}$$ 
and it contributes to the long/short margin when $s_k$ is positive/negative.

The openning position contributes the long/short margin when $s$ is positive/negative and compensate for the other one. Specifically, its pre-scale initial margin is 
$$PIM_p = s\times\max\{I,|r_m|\}$$

The pre-scale \textbf{long} margin is:
\begin{equation}\label{pim+}
    PIM^+ = \begin{cases}
        PIM_p + \sum_{k\leq bid}PIM_k & \text{ if } s + \sum_{k\leq bid}s_k > 0 \\
        0 & \text{ otherwise}
    \end{cases}
\end{equation}
    

We allow users to "take profit" here on their short position without maintain the "take profit" orders, which is when $s + \sum_{k\leq bid}s_k > 0$.

The pre-scale \textbf{short} margin is:
\begin{equation}\label{pim-}
    PIM^- = \begin{cases}
        -PIM_p + \sum_{k\geq ask}PIM_k & \text{ if } s + \sum_{k\geq ask}s_k < 0 \\
        0 & \text{ otherwise}
    \end{cases}
\end{equation}

Similar to the previous case, we allow users to "take profit" here on their long position without maintain the "take profit" orders, which is when $s + \sum_{k\leq bid}s_k < 0$.

The pre-scale initial margin of user is
$$
PIM = \max\{PIM^+,PIM^-\}
$$

The final initial margin of user is the scaled (and capped) version, where we take into account time to maturity:
$$
IM = k_{IM}\times\max\{t,t_{th}\}\times PIM
$$

User's position could be marked-to-market to determine his collateral:
\begin{equation}
    V = c + s\times r_m\times t
\end{equation}

Actions leading to potential higher leverage include: Open new order, withdraw collateral. These action can only be executed when user has enough collateral to cover the initial margin after the action:
\begin{definition}[Strong Margin Check]\label{strong-margin-check}
We allow user to do any action as long as after all action:
    $$
V \geq IM
$$
\end{definition}

One exemption is for market orders that reduce leverage and healthy for the system. This would be mention in \ref{weak-margin}.


\subsection{Liquidation}
User $i$'s state is liquidatable when:
$$
V <  MM := |k_{MM}\times \max\{t,t_{th}\}\times s\times\max\{{I,|r_m|}\}|
$$

\section{Basic Results}

\begin{prop}
    At any time $t$, for each market and for any user, the following is true:
    \begin{align}
    \sum_{i=1}^N s^{(i)}(t) = 0\\
        bid(t) < ask(t)\\
        IM(t) \geq \frac{k_{IM}}{k_{MM}} MM(t)
    \end{align}
\end{prop}
\begin{proof}
    (3.1) and (3.2) are obvious. For (3.3), we consider case by case:
    \begin{itemize}
        \item If $s \geq 0$ then $s + \sum_{k\leq bid}s_k >0$ since $s_k \geq 0$ for any $k\leq bid$. Therefore, from \eqref{pim+} we obtain:
        \begin{align*}
            IM &\geq k_{IM}\times \max\{t,t_{th}\}\times PIM^+\\
            &\geq k_{IM}\times \max\{t,t_{th}\}\times PIM_p\\
            &= k_{IM}\times \max\{t,t_{th}\}\times s\times\max\{I,|r_m|\}\\
            &=\frac{k_{IM}}{k_{MM}} k_{MM}\times \max\{t,t_{th}\}\times s\times\max\{I,|r_m|\} = \frac{k_{IM}}{k_{MM}}MM
        \end{align*}

        \item Otherwise, if $s < 0$, then $s + \sum_{k\geq ask}s_k <0$ since $s_k \leq 0$ for any $k\geq ask$. Similarly, we obtain from \eqref{pim-}:
        \begin{align*}
            IM &\geq k_{IM}\times \max\{t,t_{th}\}\times PIM^-\\
            &\geq k_{IM}\times \max\{t,t_{th}\}\times (-PIM_p)\\
            &= k_{IM}\times \max\{t,t_{th}\}\times |s|\times\max\{I,|r_m|\}\\
            &=\frac{k_{IM}}{k_{MM}} k_{MM}\times \max\{t,t_{th}\}\times |s|\times\max\{I,|r_m|\} = \frac{k_{IM}}{k_{MM}} MM
        \end{align*}
    \end{itemize}
\end{proof}
\begin{remark}
    Consider multiple markets, from (3.3) we have: 
    $$
    IM > \min_i(\frac{k_{IM}(i)}{k_{MM}(i)}) MM
    $$
    where the minimum is taken among markets. 
\end{remark}

\begin{remark}
    Note that $\frac{k_{IM}}{k_{MM}}$ is typically set to be $2$, e.g. $k_{IM} = 0.5, k_{MM} = 0.25$. Therefore, if user successfully places a new order then he can not be liquidated immediately because $$V\geq IM\geq \frac{k_{IM}}{k_{MM}} MM > MM$$.

    However, when an existing order of a user is filled, that user might become immediately liquidatable or bad debt. This is because after an order is filled, total value $V$ might change because of 1. Change in unrealized PnL of the current position and 2. The difference between the mark price and the limit price. Therefore, it might be no longer sufficient to meet the initial margin requirement at that point, i.e. $V < IM$. The user's state can jump from \textbf{healthy} ($V \geq MM$) to \textbf{unhealthy} ($V < MM$) or even \textbf{bad debt} ($V < 0$) after the order is filled. This would be discussed in the following sections.
\end{remark}
\section{Normal Market Condition}

In this section, we consider the case where the mark price follows the last traded price closely and (WLOG) the rate is always positive.

\begin{prop}
    There exists $R_0$ such that at when $r_m = R_0$, the user has exactly enough margin for the Initial Margin Requirement, i.e. $V = IM$. 
    
    Moreover, consider a user that has a long position, i.e. $s \geq 0$, then the user cannot be \textbf{liquidatable after order filled} if their orders' limit price is larger than $\max\{R_0\times (1-\frac{(k_{IM} - k_{MM})(s+s')}{(1-k_{MM})s + (k_{IM} - k_{MM})s'}), I\}$, where $s'$ is total order size of limit long orders.

    Otherwise, if a user has a short position, i.e. $s \leq 0$, then the user cannot be \textbf{liquidatable after order filled} if their orders' limit price is smaller than $R_0\times (1+\frac{(k_{IM} - k_{MM})(s+s')}{(1+k_{MM})s - (k_{IM} - k_{MM})s'})$, where $s'$ is total order size of limit long orders.
\end{prop}
\begin{proof} 
    \textit{Case 1: Long position.} Let $r_{m'}$ denotes the smallest price among the user's limit long orders' price. Assume $r_m = R_0$, then
    $$
    V =  IM\geq k_{IM}\times s\times t\times r_m + k_{IM}\times s_{m'}\times t\times r_{m'}
    $$ 

    Suppose that the order at $r_{m'} = r_m - \Delta > I$ is filled and user is liquidatable immediately after that. By assumption, the mark price follows closely the last traded price, meaning user doesn't satisfy the Maintenance Margin Requirement at $r_{m'}$. Therefore:
    \begin{align*} 
         &k_{MM}\times t\times (s+s')\times (r_m-\Delta)> V - s\times t\times\Delta\\
         &\geq k_{IM}\times s\times t\times r_m + k_{IM}\times s_{m'}\times t\times r_{m'} - s\times t\times\Delta\\
        &\iff (k_{IM}-k_{MM})\times(s + s_{m'})\times t\times r_m \leq \Delta\times(s\times t - k_{MM}\times t \times (s + s_{m'}) + k_{IM}\times s_{m'}\times t)\\
        &\iff \frac{\Delta}{r_m} \geq \frac{(k_{IM} - k_{MM})(s+s')}{(1-k_{MM})s + (k_{IM} - k_{MM})s'}
    \end{align*}
    which is contradiction to the assumption that $r_m' > r_m\times (1-\frac{(k_{IM} - k_{MM})(s+s')}{(1-k_{MM})s + (k_{IM} - k_{MM})s'})$.

    \textit{Case 2: Short position.} Let $r_{m'}$ denotes the highest price among the user's limit short orders' price. Assume $r_m = R_0$, then
    $$
    V =  IM\geq -k_{IM}\times s\times t\times r_m - k_{IM}\times s_{m'}\times t\times r_{m'}
    $$ 

    Similarly to the previous section, user is liquidatable immediately after the order at $r_{m'} = r_m + \Delta$ is filled means:
    \begin{align*} 
         &-k_{MM}\times t\times (s+s')\times (r_m+\Delta)> V + s\times t\times\Delta\\
         &\geq -k_{IM}\times s\times t\times r_m -k_{IM}\times s_{m'}\times t\times r_{m'} + s\times t\times\Delta\\
        &\iff (k_{IM}-k_{MM})\times(s + s_{m'})\times t\times r_m > \Delta\times(s\times t + k_{MM}\times t \times (s + s_{m'}) - k_{IM}\times s_{m'}\times t)\\
        &\iff \frac{\Delta}{r_m} > \frac{(k_{IM} - k_{MM})(s+s')}{(1+k_{MM})s - (k_{IM} - k_{MM})s'}
    \end{align*}
    which is contradiction to the assumption that $r_m' < r_m\times (1+\frac{(k_{IM} - k_{MM})(s+s')}{(1+k_{MM})s - (k_{IM} - k_{MM})s'})$.
\end{proof}

\begin{remark}
    When user max long/short, that is, he uses up all of his margin for opening a long/short order, then the current mark price is $R_0$ of his current state.
\end{remark}

\begin{remark}
    When $s' = 0$, the formula equals to the liquidation price: $R_0\times(1-\frac{k_{IM}-k_{MM}}{1-k_{MM}})$ for longing position and $R_0\times(1+\frac{k_{IM}-k_{MM}}{1+k_{MM}})$
\end{remark}

\begin{remark}
    Let $k_{IM}=0.5,k_{MM}=0.25$ and $r_m$ is the current mark price. Suppose user has an opening position of size $s>0$. If he opens maximum long order size and the order size is $\af\times s$, then he would be liquidatable after the order is filled if the limit price is lower than $r_m\times(1- \frac{1+\af}{3+\af})$.
\end{remark}

\begin{remark}
    We give an example to illustrate why this threshold is tight. Let $k_{IM}=0.5,k_{MM}=0.25$ and $r_m=R_0 = 20\%$ is the current mark price. 
\end{remark}

\section{Extreme market condition}

Under extreme market condition, the assumption that $r_m$ is closed to the last traded price might not be correct. In this section, we consider the case where orders with limit price away from $r_m$ is filled. Once again, we only consider the case where the rate is always positive. 

Specifically, we look into the events where user places a limit order. Let $t_1,t_2$ denote the time where he places the order and when his order is filled. We assume $1 > k_{IM} > k_{MM}$ in the following results. 
\begin{prop}
    Assume that $I = 0$. Assume that at $t_2-$, user's state satisfies the IM requirement. If only the long orders at price smaller than $r_m\times\frac{2}{2-k_{IM}}$ and short order at price larger than $r_m\times\frac{2}{2+k_{IM}}$ can be filled then user cannot be in bad debt at $t_2+$.
\end{prop}
\begin{proof}
Let $(s_{m'},r_m)'$ denote his highest price long order. By assumption that user's state satisfies the IM requirement at $t_2-$:
    \begin{align}
    V >  IM&\geq |k_{IM}\times s\times t\times r_m + k_{IM}\times s_{m'}\times t\times r_{m'}|\\
    \text{and } V >  IM&\geq k_{IM}\times |s|\times t\times r_m
    \end{align}

After filled:
$$
V^+ = c + s\times t\times r_m - s_{m'}\times t\times (r_{m'}-r_m)
$$

% Suppose user is liquidatable after filled:
% \begin{align*}
% &MM = k_{MM}\times (s+s_{m'})\times t\times r_m \\
% &> c + s\times t\times r_m - s_{m'}\times t\times (r_{m'}-r_m) \\
% &\geq k_{IM}\times s_{m'}\times t\times r_{m'} - s\times t\times r_m\times (1-k_{IM}) +  s\times t\times r_m - s_{m'}\times t\times (r_{m'}-r_m)\\
% \RA &\frac{\Delta}{r_m}s_{m'} > (s+s_{m'}) \frac{k_{IM} - k_{MM}}{1 - k_{IM}}
% \end{align*}

Suppose user is in bad debt after filled, from (5.1):
\begin{align*}
&0> c + s\times t\times r_m - s_{m'}\times t\times (r_{m'}-r_m) \\
&\geq s\times t\times r_m\times k_{IM} + k_{IM}\times s_{m'}\times t\times r_{m'} - s_{m'}\times t\times (r_{m'}-r_m)\\
\RA &\frac{r_{m'}}{r_m} > \frac{sk_{IM} + s_{m'}}{s_{m'}(1-k_{IM})}
\end{align*}

On the other hand, from (5.2):
\begin{align*}
&0> c + s\times t\times r_m - s_{m'}\times t\times (r_{m'}-r_m) \\
&\geq |s|\times t\times r_m\times k_{IM} - s_{m'}\times t\times (r_{m'}-r_m)\\
\RA &\frac{r_{m'}}{r_m} > 1 + k_{IM}\frac{|s|}{s_{m'}}
\end{align*}

Therefore: 
$$
\frac{r_{m'}}{r_m} > \max_{s_{m'}>0}\{\frac{sk_{IM} + s_{m'}}{s_{m'}(1-k_{IM})},1 + k_{IM}\frac{|s|}{s_{m'}}\} \geq \frac{2}{2-k_{IM}} \text{ (contradiction)}
$$

The last equality happens when $\frac{s_{m'}}{s} = k_{IM}-2$.

Similarly, if we let $(s_{m'},r_m)'$ denote the lowest price ask order, then we would have: 
$$
\frac{r_{m'}}{r_m} < \min_{s_{m'}<0}\{\frac{-sk_{IM} + s_{m'}}{s_{m'}(1+k_{IM})},1 + k_{IM}\frac{|s|}{s_{m'}}\} \leq \frac{2}{2+k_{IM}}\text{ (contradiction)}
$$

The last equality happens when $\frac{s_{m'}}{s} = -(k_{IM}+2)$.

\end{proof}
\begin{remark}\label{6.1}
    We also show that this bound is tight by an example. Assume that $k_{IM} = 0.5$, $t=1$ and the current mark rate is $r_m = 15\%$. 
    
    Case 1: A user has a short position $s=-10$. He places a long order at $r_{m'} = r_m\frac{2}{2-k_{IM}} = 20\%$ of size $s_{m'} = s(k_{IM}-2) = 15$. 
    
    The Initial Margin of this user in this case is 
    $$
    IM = \min\{k_{IM}\times |s|\times t\times r_m,|k_{IM}\times (s\times t\times r_m+s_{m'}\times t\times r_{m'})\} = 0.75
    $$
    
    He used all of his collateral for this order, meaning $V = IM = 0.75$, or $c = 0.75 -s\times t \times r_m = 2.25$. After this order is filled, his new total value is $V - s_{m'}\times(r_{m'} - r_m) = 0.75 - 15(20\%-15\%)=0$, making him in bad debt immediately.

    Case 2: A user has a long position $s=10$. He places a short order at $r_{m'} = r_m\frac{2}{2+k_{IM}} = 12\%$ of size $s_{m'} = -s(k_{IM}+2) = -25$. 
    
    The Initial Margin of this user in this case is 
    $$
    IM = \min\{k_{IM}\times |s|\times t\times r_m,|k_{IM}\times (s\times t\times r_m+s_{m'}\times t\times r_{m'})\} = 0.75
    $$
    
    Similarly, consider the case where he used all of his collateral for this order, i.e. $V = 0.75$ or $c = 0.75 -s\times t \times r_m = -0.75$. After this order is filled, his new total value is $V - s_{m'}\times(r_{m'} - r_m) = 0.75 + 25(12\%-15\%)=0$, making him in bad debt immediately.
\end{remark}

We extend Proposition 6.1 into the following setting, where we tighten the condition on the order's price volatility and take into account $I$ when calculating margins.

\begin{prop}\label{core-limit-range}
    Consider one market. Assume that at $t_2-$, user's state satisfies the IM requirement. There exists $f^u(r_m),f^l(r_m)$ such that: If only the long orders at price smaller than $f^u(r_m;k_{MM},k_{IM})$ and short orders at price larger than $f^l(r_m;k_{MM},k_{IM})$ can be filled then user's health is always larger than 1 immediately after that limit order is filled (he is not liquidatable).  
\end{prop}
\begin{proof}
    
We apply the same setting as before, where $(s_{m'},r_{m'})$ denote his highest price long order to be filled. 
    \begin{align}
    V >  IM&\geq |k_{IM}\times s\times t\times r_m^+ + k_{IM}\times s_{m'}\times t\times r_{m'}^+|\\
    \text{and } V >  IM&\geq k_{IM}\times |s|\times t\times r_m^+
    \end{align}
where $r^+$ denotes $\max(I, r)$.

After filled:
$$
V^+ = c + s\times t\times r_m - s_{m'}\times t\times (r_{m'}-r_m)
$$

Suppose user's health is smaller than 1 after the order is filled, from (6.3):
\begin{align*}
& k_{MM}\times |s+s'|\times t\times r_m^+> c + s\times t\times r_m - s_{m'}\times t\times (r_{m'}-r_m) \\
&\geq k_{IM}(s\times t\times r_m^+ + s_{m'}\times t\times r_{m'}^+)  - s_{m'}\times t\times (r_{m'}-r_m)\\
\RA g(r_{m'}) &\geq r_m + r_m^+(k_{IM}\times \frac{s}{s_{m'}} -k_{MM}\times|1+\frac{s}{s_{m'}}|) =: f_1(x;r_m)|_{x=\frac{s}{s_{m'}}}
\end{align*}
where $g(x) := x - k_{IM}\times x^+$.

On the other hand, from (6.4):
\begin{align*}
& k_{MM}\times |s+s'|\times t\times r_m^+ > c + s\times t\times r_m - s_{m'}\times t\times (r_{m'}-r_m) \\
&\geq |s|\times t\times r_m^+\times k_{IM} - s_{m'}\times t\times (r_{m'}-r_m)\\
\RA & r_{m'} > r_m + r_m^+(k_{IM}|\frac{s}{s_{m'}}| -  k_{MM}\times|1+\frac{s}{s'}|) =: f_2(x;r_m)|_{x=\frac{s}{s_{m'}}}
\end{align*}

Since $g$ is an increasing function:
$$
r_{m'} \geq \max(g^{-1}(f_1(x;r_m)), f_2(x;r_m)) \geq f^u(r_m)
$$
where $f^u(r_m) := \min_x\{\max(g^{-1}(f_1(x;r_m)), f_2(x;r_m))\}$. 

Since $r_{m'} > r_m(t_2-)$, it's also true that $r_{m'} > r_m(t_1+)$ because during $(t_1+,t_2-)$, there is no transaction between $r_{m'}$ and $r_m(t_1+)$, and hence $r_{m'} > r_m(t_2-) \geq r_m(t_1+)$. This contradicts our assumption about the upper bound of the long orders.

Similarly, consider $(s_{m'},r_{m'})'$ which is the lowest bid order. Supposer user's health is smaller than 1 after this order is filled:
\begin{align*}
    h(r_m) := r_{m'} + k_{IM} r_{m'}^+ &\leq r_m - r_m^+(k_{IM}x - k_{MM}|1+x|) := f_3(x;r_m)|_{x=\frac{s}{s_{m'}}}\\
    r_{m'} &\leq r_m - r_m^+(k_{IM}|x| - k_{MM}|1+x|) := f_4(x;r_m)|_{x=\frac{s}{s_{m'}}}\\
    \RA r_{m'} &\leq \min(h^{-1}(f_3(x;r_m)), f_4(x;r_m)) \leq f^l(r_m)
\end{align*}
where $f^l(r_m) := \max_x\{\min(h^{-1}(f_3(x;r_m)), f_4(x;r_m))\}$. 

Since $r_{m'} < r_m(t_2-)$, it's also true that $r_{m'} < r_m(t_1+)$ because during $(t_1+,t_2-)$, there is no transaction between $r_{m'}$ and $r_m(t_1+)$, and hence $r_{m'} < r_m(t_2-) \leq r_m(t_1+)$. We have the contradiction here by the assumption on the lower bound of the short orders.

\end{proof}
\begin{remark}
    We derive $f^u, f^l$ in Appendix \ref{limit-range}. Similar to Remark \ref{6.1}, we can show that this bound is tight. 
\end{remark}

\begin{coro}
    Let $\af, \bt\in[0,1]$ such that $1 > \bt k_{IM} > \af k_{MM}$. If we assume that at $t_2-$, user's state satisfies $V > \bt IM$, and the limit price range $[f^u(r_m;\af k_{MM}, \bt k_{IM}), f^l(r_m;\af k_{MM}, \bt k_{IM})]$ is applied, then we can conclude that after the order is filled at $t_2$, user's health ratio cannot drop below $\af$.
\end{coro}

\begin{coro}
    Consider multiple market. Assume that at $t_2-$, user's state satisfies the IM requirement. If \textbf{on every market}, the limit price range $[\min f^u(r_m; k_{MM}, k_{IM}), \max f^l(r_m; k_{MM}, k_{IM})]$ then user's health is always larger than 1 immediately after that limit orders are filled (he is not liquidatable).
\end{coro}

\begin{proof}
    Suppose there are $M$ markets, let $(s'_i, r_{m^i})$ denote the order to be filled on market $i$ and $s_i$ denote the current position size on market $i$. Since user satisfies the IM requirement at $t_2-$: 
    $$
    c + \sum_i s_i\times t\times r^i_m > \sum_i IM(i)
    $$

    Suppose the user's health is smaller than 1 after the order is filled: 
    \begin{align*}
    \sum_i MM(i) &> c + \sum_i s_i\times t\times r^i_m - \sum_i s_i'\times t\times (r_{m^i}-r^i_m)\\
    &> \sum_i IM(i) - \sum_i s_i'\times t\times (r_{m^i}-r^i_m)
    \end{align*}

    Therefore, there exists at least one market such that:
    $$
    MM(i) > IM(i) - s_i'\times t\times (r_{m^i}-r^i_m)
    $$

    Follow the arguments in Proposition 6.2, we can show that it either violates the upper bound when $s_i' > 0$
    $$
    r^i_m > f^u_i(r_m) \geq \min f^u(r_m; k_{MM}, k_{IM})
    $$  
    or the lower bound when $s_i' < 0$
    $$r^i_m < f^l_i(r_m) \leq \max f^u(r_m; k_{MM}, k_{IM})$$

    We conclude the proof.
    
\end{proof}

Combine Corollary 6.1 and Corollary 6.2, we have the following result to define our limit price range.

\begin{coro}\label{5.3}
    Let $\af, \bt\in[0,1]$ such that $1>\bt k_{IM} > \af k_{MM}$. Consider multiple market. Assume that at $t_2-$, user's state satisfies $V > \bt IM$. If \textbf{on every market}, only the long orders at price smaller than $f^u(r_m;\af k_{MM}, \bt k_{IM})$ and short order at price larger than $f^l(r_m;\af  k_{MM}, \bt k_{IM})$ can be filled then user's health is always larger than $\af$ immediately after that limit order is filled.
\end{coro}

With this result, we give a definition of Limit Order Range.
\begin{definition}[Limit Order Range]
    Given market price $r_m$, users are only allowed to place limit long orders at price smaller than $f^u_s(r_m;H_c k_{MM}, k_{IM})$ and short orders at price larger than $f^l_s(r_m;H_c k_{MM}, k_{IM})$. 
    The formula for this bound is in Appendix \ref{limit-range}. 
\end{definition}

\begin{remark}
    In some scenarios, we might set $k_{IM} > 1$. In order to guarantee our result, we adapt the price ranges $f^u_s(r_m; H_c k_{MM}, \min(0.9, k_{IM}))$ and $f^l_s(r_m; H_c k_{MM}, \min(0.9, k_{IM}))$ rather than extending these bounds according to $k_{IM} > 1$. Corollary 6.3 still follows when we set $\beta k_{IM} = 0.9$. That is, we still ensure that if a user satisfies the IM requirement at $t_2-$, i.e., $V > IM \geq \beta IM$, their health after these orders are filled is greater than $\alpha$.
\end{remark}

\section{Automatic Bots}
We extend the definition of user's health here to take into account user's orders as well.
\begin{definition}[General Health]\label{H-after}
    For any set of open orders $O$, let $H(O)$ denote user health AFTER all these orders are filled. The current user's health is $H := H(\emptyset)$.
    
    $H_w$ is the worst possible health factor for a user if any set of open orders is filled immediately:
    $$
    H_w = \min_{O} \{H(O)\}.
    $$
\end{definition}
Although it's computationally infeasible to calculate the $H_w$ for each user, it's relatively easy to check if $H_w>C$ for any constant $C$. 
\begin{assumption}[Normal Bots]\label{bots} We have the following bots:
    
    (i) Bot 1 is allowed to deleverage user when $H < H_d$.
    
    (ii) Bot 2 is allowed to cancel user's open order when $H_w < H_c$. 
    
\end{assumption}

\begin{remark}
    The condition of bot 2 is weaker than bot 1 since $H < H_d \RA H_w\leq H < H_d < H_c$.
\end{remark}
\begin{remark}
    By Corollary \ref{5.3}, if user satisfies IM requirement then $H_w > H_c$. 
\end{remark}

\section{Closing Order margin check}\label{weak-margin-check}
We allow users with low margin (insufficient initial margin) to close their order in a healthy way. 
\begin{definition}[Closing Order Margin Check]
    We allow close orders of size $s'$ at $r'$ such that:
    \begin{equation}\label{weak-margin}
    \left|\frac{(r'-r)}{ r^+}\right|\leq k_{MM}\times H_f 
    \end{equation}
    and $0< s + s' < s$ or $0 > s+s' > s$.

    These orders are refereed to as ``weak order".
\end{definition}

\begin{remark}
    Under normal market condition, these conditions always satisfy if user market close since $r\approx r'$.
\end{remark}

\begin{remark}
    Note that $s'(r'-r)$ and $k_{MM}\times|s'|\times r^+$ is the change in user's total value and maintenance margin after the order is filled.
\end{remark}

This implies that after the order is filled, the health ratio of user improves. We show this in the following proposition.

\begin{assumption}\label{liveness}
    Whenever there is a highest bid order with price higher than last traded price or a lowest ask order with price lower than last traded price, the mark price would be updated as if that order price is the last traded price.
\end{assumption}
\begin{prop}\label{8.1}
    Consider multiple market. Let $t_2$ be the time where the order passing a weak margin check is filled. Assume that at $t_2-$ (i.e. right before the order is filled), health ratio of user is larger than $H_f$, then at $t_2+$ (i.e. right after the order is filled), user's health increases, and hence larger than $H_f$.
\end{prop}
\begin{proof}
    Let $t_1$ be the time of order placement. Let health of the user at time $t_2-$ be $\frac{V}{MM}>H_f$. Let the mark rate at $t_2$ be $r(t_2)$. WLOG, assume user is having a short position and placed a long closing order ($s(t_1) < 0$ and  $s' > 0$). 

    Assuming $H(t_2-) = \frac{V}{MM},$ health of user at $t_2+$ is
    $$
    H(t_2+) = \frac{V-s'(r'-r(t_2))}{MM-k_{MM}\times|s'|\times r(t_2)^+}.
    $$
    
    We have the condition \eqref{weak-margin}:
    $$
    \left|\frac{s'(r'-r(t_1))}{k_{MM}\times s'\times r(t_1)^+}\right| \leq H_f
    $$

    If $r' \leq r(t_2)$, then obviously $H(t_2+) > H(t_2-)$, which implies that user's health improves after the order is filled.

    Otherwise, we consider the case $r(t_2) < r'$. 

    If $r(t_1) > r'$: Note that the order is not filled during $(t_1+, t_2-)$, therefore, all the orders filled during this duration are of price at least $r'$. If there is at least one other order filled during this duration, then mark rate is always updated above $r'$, i.e. $r(t_2) \geq r'$ (contradiction). Similarly, if there is no order filled during this duration, and last traded price is of price at least $r'$, then $r(t_2) \geq r'$ (contradiction). Otherwise, if there is no order filled during this duration, and last traded price is of price less than $r'$, then at $t_1$, mark rate would be updated once again at $r'$ since this bid order price is higher than the last traded price by assumption \label{liveness}. Therefore, we have $r(t_2) \geq r'$ (contradiction).
    
    If $r(t_1) > r'$,  If $r(t_1) \leq r(t_2) < r'$, then:
    $$
    \frac{s'(r'-r(t_2))}{k_{MM}\times s'\times r(t_2)^+} \leq \frac{s'(r'-r(t_1))}{k_{MM}\times s'\times r(t_1)^+} \leq H_f
    $$
    which, together with $\frac{V}{MM} \geq H_f$, implies
    $$
    H(t_2+) \geq H(t_2-)
    $$
    
    
\end{proof}
\begin{prop}\label{8.2}
    Consider multiple market. Placing a new order passing Weak Margin Check only improves user's $H_w$, assuming that $H_w > H_f$. 
\end{prop}
\begin{proof}
    Suppose $H_w-$ at $t_2-$ (the moment right before order placement of $o$) is:
    $$
    H_w- = H(O)
    $$
    
    Consider $H_w+$ at $t_2+$:
    $$
    H_w+ = H(O')
    $$
    
    If the new order $o$ is not in $O'$, then $O' = O$, i.e. $H_w+=H_w-$.

    Otherwise, if $o\in O'$, then $H(O'\setminus o) \geq H(O) > H_f$. 

    Similar to Proposition \ref{8.1}, we have $H_w+ = H(O') \geq H(O'\setminus o) \geq H_w-$.

    We conclude the proof here.
\end{proof}



\section{Mark Price Volatility}\label{volatility}
Besides sudden changes in a user's health ratio from open orders or closed positions, the mark price of the market also affects their health ratio.

To prevent abrupt and rapid drops in a user's health ratio that could lead to liquidations or global deleveraging occurring too late to prevent bad debt, we limit the rate of change in the mark price. This gradual adjustment allows the system sufficient time to react.

Specifically, we don't allow the transaction to be significantly far away from the current mark price. 

Throughout this section, we use $C = 5\text{ min}$ as the unit of times, so $t=1$ corresponds to 5 minutes.
\begin{definition}[General Order Restriction/Max Rate Deviation]
Let $r_n$ denote the new limit price and $r_m$ be the current mark price. We require:
$$
|r_n - r_m| < k_{MD}\times \max(|r_m|,I) 
$$
where $k_{MD}$ is a constant, chosen to be less than $k_{MM}$:
\[
k_{MD} \leq k_{MM}
\]
    
\end{definition}

Recall the calculation for mark price when the last trade happens on rate $r$ at time $t_1$:
$$
r_m(t_2) = r\times(t_2-t_1) + r_m(t_1)\times (1-(t_2-t_1))
$$

From this, we obtain:
\begin{equation}\label{8.1}
    |\frac{\de r_m}{\de t}| \leq k_{MD} |r_m|^+.
\end{equation}

% \subsection{Upperbound} If $r(t_1) \geq I$ then $r(t_2) < r(t_1)\exp(k(t_2-t_1))$. Otherwise, consider the case $r(t_1) < I$. Let $t^*$ denote the time such that $r(t_1)+kI(t^*-t_1) = I$. 

% If $t_2 < t^*$ then $r(t_2) < r(t_1)+kI(t_2-t_1)$. Otherwise, if $t_2 > t^*$ then $r(t_2) < I\exp(k(t_2-t^*))$.

% \subsection{Lowerbound} If $r(t_1) < I$ then $r(t_2) < r(t_1)-kI(t_2-t_1)$. Otherwise, consider the case $r(t_1) > I$. Let $t^*$ denote the time such that $r(t_1)\exp(-k(t^*-t_1)) = I$. 

% If $t_2 < t^*$ then $r(t_2) < r(t_1)\exp(-k(t_2-t_1))$. Otherwise, if $t_2 > t^*$ then $r(t_2) < I-kI(t_2-t^*)$.

% \begin{remark}
%     Suppose we take $k=k_{IM} = 0.5$, $C = 5$ minutes and $I = 10\%$. After 30 seconds, mark price increases/decreases by at most $\approx 5\%$ of the current price when price is larger than $I$, or absolute $0.5\%$ when price is smaller than $I$.
% \end{remark}

\subsection{Health Volatility}\label{H-volatility}
\begin{definition}
    We define $r^u(r, dt)$ as the upper bound of $r(t+dt)$ given that $r(t) = r$. Similarly, $r_l(r, dt)$ is the corresponding lower bound. We define $Range(r,t) := [r_l(r,t), r^u(r,t)]$. 
\end{definition}

\begin{definition}[Worst Health Jump]
    Define:
        $$
        f(r;c,s) := \frac{c + sr}{k_{MM}|s|r^+}
        $$
    We define the worst health jump $J(H,t)$ as following:
    \begin{equation}
        J(H,t) = \max_{c,s,r: f(r;c,s) = H, r'\in Range(r,t)} \left(f(r;c,s) - f(r';c,s) \right)
    \end{equation}
\end{definition}

\begin{prop}
    We have the following properties of $J$:

    1. For any $H>0$:
    \begin{equation}
        J(H,0) = 0
    \end{equation}

    2. For any $t_1 \leq t_2$, $H>0$:
    \begin{equation}
        J(H,t_1) \leq J(H,t_2)
    \end{equation}
    
    % 3. Define $l(H,t) := H - D(H, t)$. For any $0 < H_1 \leq H_2$ and $t>0$, we have:
    % \begin{equation}\label{8.4}
    % l(H_1, t) \leq l(H_2,t)
    % \end{equation}

    3. For any $t>0$ and $H\in((1+k_{MM}) t,1)$, we have:
    \begin{equation}\label{8.5}
        J(H,t) \leq \frac{k_{MD}}{k_{MM}}(1 + k_{MM}) t =: D(t)
    \end{equation}

    % 4. For any $H'_1 \geq l(H_1, t_1)$, $H'_2 \geq l(H_2, t_2)$ and $H'_1 \leq H_2$, we have:
    % \begin{equation}\label{8.5}
    %     H_1 - H'_2 \leq D(H_1, t_1+t_2)
    % \end{equation}
\end{prop}
\begin{proof}
    % We prove \eqref{8.4}. We write $l(H) = l(H,t)$ in the context of this proof. By definition of $D$, we have:
    % \begin{equation}
    %     l(H) = \min_{c, s, r'\in Range(r,t), f(r;c,s) = H} f(r'; c,s)
    % \end{equation}
    
    % Suppose in contrast, we have $H_2 \geq H_1 \geq l(H_1) > l(H_2)$ with the corresponding argmin $(c_1, s_1, r_1, r_1')$ and $(c_2, s_2, r_2, r_2')$. 

    % Consider the path from $r_2$ to $r_2'$. Since $f(r; c_2, s_2) := \frac{c_2 + s_2r}{b|s_2|r^+}$ is a continuous function w.r.t. $r$, there exists $r_3, r_4$ in this path such that $f(r_3; c_2, s_2) = H_1$ and $f(r_3; c_2, s_2) = l(H_1)$. Obviously, since the mark rate can go from $r_3$ to $r_2'$ in the time duration $t$, we have $r_2'\in Range(r_3,t)$. Therefore, by definition: 
    % $$
    % l(H_1) = \min_{c, s, r'\in Range(r), f(r;c,s) = H_1} f(r'; c,s) \leq f(r_2'; c_2, s_2) = l(H_2),
    % $$
    % which is a contradiction. 

    We prove \eqref{8.5}. We only need to show that:
    \begin{equation}
        \left|\frac{\partial f(r(t);c,s)}{\partial t} \right| \leq  1 + k_{MM}
    \end{equation}
    for any $c,s, r(t)$ such that $f(r(t);c,s) \in((1+k_{MM}) t,1)$.

    \textbf{Case 1}: $r > I$. We show that
    \begin{equation}
        \left|\frac{\partial f(r;c,s)}{\partial r} \right| \leq  \frac{k_{MM} + 1}{k_{MM}r}
    \end{equation}

    WLOG, assume $|s| = 1$.
    
    \textbf{Sub-case 1.1}: $s = 1$. The condition $0 < f(r(t);c,s) < 1$ implies $-1 < \frac{c}{r} < -(1-k_{MM})$. 

    By definition:
    \begin{align}
        \left|\frac{\partial f(r;c,s)}{\partial r} \right| = \left|\frac{c}{k_{MM}r^2}\right| < \frac{1}{k_{MM}r}\ < \frac{k_{MM} + 1}{k_{MM}r}
    \end{align}

    \textbf{Sub-case 1.2}: $s = -1$. The condition $0 < f(r(t);c,s) < 1$ implies $1 < \frac{c}{r} < (k_{MM} + 1)$. 

    By definition:
    \begin{align}
        \left|\frac{\partial f(r;c,s)}{\partial r} \right| = \left|\frac{c}{k_{MM}r^2}\right| < \frac{k_{MM} + 1}{k_{MM}r}
    \end{align}

    By chain rule, using \eqref{8.1}, we have:
    \begin{equation}
        \left|\frac{\partial f(r(t);c,s)}{\partial t} \right| = \left|\frac{\partial f(r;c,s)}{\partial r} \right| \left|\frac{dr}{t} \right| \leq \frac{k_{MM} + 1}{k_{MM}r} \times k_{MD} r = \frac{k_{MD}}{k_{MM}}(1+k_{MM})
    \end{equation}

    \textbf{Case 2}: $r < I$. 

    By definition:
    \begin{align}
        \left|\frac{\partial f(r;c,s)}{\partial r} \right| = \frac{1}{k_{MM}I} 
    \end{align}

    By chain rule, using \eqref{8.1}, we have:
    \begin{equation}
        \left|\frac{\partial f(r(t);c,s)}{\partial t} \right| = \left|\frac{\partial f(r;c,s)}{\partial r} \right| \left|\frac{dr}{dt} \right| \leq \frac{1}{k_{MM}I}  \times k_{MD} I < \frac{k_{MD}}{k_{MM}}(1 + k_{MM})
    \end{equation}

    We conclude the proof here.

    % We prove \eqref{8.5}:
    % \begin{align}
    %     H_1 - H'_2 &= H_1 - H_2 + H_2 - H'_2\\
    %     &\leq H_1 - H_2 + D(H_2,t_2)\\
    %     &\leq H_1 - H_1' + D(H_1',t_2)\\ 
    %     &= H_1 - H'_1 + H'_1 - H_3\\
    %     &= H_1 - H_3
    % \end{align}
    % where some $H_3 = f($
    
\end{proof}

\begin{remark}\label{P-bound}1
    Consider multiple markets. We also have $J(H,t) \leq D(t) = \max_i \{\frac{k_{MD}(i)}{k_{MM}(i)}(1+k_{MM}(i))t\}$ where $k_{MM}(i)$ denote the $k_{MM}$ from market $i$.
\end{remark}

\subsection{Worst Health Volatility}\label{Hw-P-bound}
Our checks take into account of $H_w$ in \ref{H-after}. Assume that $H_w = H(O)$:
$$
H_w = \frac{c - \sum_{o\in O}s_o\times r'_o + \sum_{o\in O}s_o\times r_o}{\sum_{o\in O} b_o\times s_o\times\max(r_o, I_o)}
$$
where $s_o, r_o, r'_o$ is the size of user AFTER order $o$ is filled, the market price and the limit price of the order in the corresponding market.

The ratio $H(O)$ behaves exactly the same as the health in the previous section as the rate changes for any fixed $O$. Therefore, for any fixed set of order $O$:
$|H'(O) - H(O)| \leq D(t)$ where $H'$ is the health after the rate changes in a duration of $t$.

Therefore:
$$
H'_w = \min_O H'(O) = H'(O^*) \geq H(O^*) - D(t) \geq \min_O(H(O) - D(t)) = H_w - D(t).
$$

We impose an assumption using this $D$ applied to the health threshold:
\begin{assumption}
    $H_c, H_d, H_f$ in this paper satisfy:
    \begin{align}
        H_c - D(30s) \geq H_d\\
        H_d - D(30s) \geq H_f
    \end{align}
\end{assumption}


\section{Safe Margin in single timestamp}
In this section, we explain why our margin system remains safe during a period without settlement. Therefore, in our analysis, we consider a fixed time $t$, and WLOG, we assume $t=1$.

We study user's health. The system is safe as long as no user has negative health.

There are two types of changes in health: (1) gradual changes due to mark price fluctuations and (2) sudden changes caused by user actions.

We limit the rate of the first kind of change by limiting the rate of change of the mark rate, as described in Section \ref{volatility}.

We visit the state changes from WHITEPAPER.

User's health in each market is determined by the tuple 
\[H(c,s,r,t) = \frac{c + s\times t \times r}{k_{MM}\times s \times t^+ \times |r|},\] 
where $c$ is the collateral, $s$ is the position size, $r$ is the mark rate, $t$ is the time until maturity.

This tuple $(c,s,r,t)$ is changed under the following "actions":
\begin{enumerate}
    \item Strong action (3.3.11-(i)): At the end of this action, we have $V > IM$. $r,t$ is not allowed to change. $c,s$ are allowed to change.
    \item Fill action (3.3.7): Some orders in \textbf{Strong} action are filled. $r,t$ is not allowed to change. $c,s$ are allowed to change.
    \item Weak action: At the end of this action, we have \textbf{Only closing orders check} satisfies. $c,s$ are allowed to change. $r,t$ is not allowed to change.
    \item Fill Weak action: Some orders in \textbf{Weak} action are filled. $r,t$ is not allowed to change. $c,s$ are allowed to change.
    \item P action: $c,s,t$ are not allowed to change. $r$ is allowed to change.
    \item T action: $s,r$ are not allowed to change, $c,t$ is changed.
\end{enumerate}

\begin{assumption}\label{fast-bot}
    It takes less than 30s for bot 1 and bot 2 to check any user.
\end{assumption}
Under this assumption, we apply the rate of change in section \ref{volatility} into our analysis.

\begin{prop}\label{safe-core-1}
    Given a time $t>60s$ such that: 

    1. There has not been any force deleverage yet (ie, no force develerage in $[0,t]$).
    
    2. There has not been any floating rate accumulation in $[t-60s, t]$.

Then $H(t-) \geq H_f$.

\end{prop}

\begin{proof}
    

We have a few types of action that affect $H = f(r,c,s)$ and $H_w$, including the special type of action related to passing of time which would affect mark rates in formulas:
\begin{itemize}
    \item \textbf{Strong} action: At the end of this action, we have $V > IM$. $r$ is not allowed to change. $c,s$ are allowed to change.
    \item \textbf{Fill Strong} action: Some orders in \textbf{Strong} action are filled. $r$ is not allowed to change. $c,s$ are allowed to change.
    \item \textbf{Weak} action: At the end of this action, we have \textbf{Only closing orders check} satisfies. $c,s$ are allowed to change. $r$ is not allowed to change.
    \item \textbf{Fill Weak} action: Some orders in \textbf{Weak} action are filled. $r$ is not allowed to change. $c,s$ are allowed to change.
    \item \textbf{P} action: $c,s$ are not allowed to change. $r$ is allowed to change.
\end{itemize}

We have the following relevant claims:
\begin{itemize}
    \item \textbf{Strong} action: At the end of this action, we have:
    \begin{equation}
        H_w(a+) \geq H_c.
    \end{equation}
    \item \textbf{Fill Strong} action: After this action, $H_w$ improves: 
    \begin{equation}
        H_w(a+) \geq H_w(a-).
    \end{equation}
    \item \textbf{Weak} action: Actions of this type do not change the ``worst order set" $O$, which is the set of order such that $H_w = H(O)$: 
    \begin{equation}
        O(a+) = O(a-)
    \end{equation}
    \item \textbf{Fill Weak} action: After this action, $H_w$ improves: \begin{equation}
        H_w(a+) \geq H_w(a-).
    \end{equation}
    \item \textbf{P} action: The rate of change of $H$ and $H_w$ caused by this action is bounded by $D(t)$ as mentioned in Remark \ref{P-bound} and \ref{Hw-P-bound}: 
\end{itemize}

We look at time $t_1 = t-30s$. Since there is no deleverage in up until time $t$, we have: 
\begin{equation}\label{10.1}
    H(t_1) > H_d
\end{equation}
Otherwise, bot 1 triggered and deleveraged user before $t_1 + 30s = t$. 

We consider $H_w$. 

\textbf{Case 1: If $H_w(t_1) > H_d$:} Consider the period from $t_1$ to $t$. 

Let $u$ be the last time of a strong action during $(t_1, t)$, and we set $u=t_1$ if there is no new strong order during this period. We have $H_w(u+) > H_c > H_d$.

Suppose all the action during $(u+, t)$ are $\{b_0, a_1, b_1, ..., a_k, b_k\}$ where $a_i$ is a action of type \textbf{Fill Strong}, \textbf{Weak} or \textbf{Fill Weak}, and $b_i$ is a \textbf{P} action, which is the change of $r$ from time between 2 actions $a_i, a_{i+1}$. We then have $H_w(a_i+) = H_w(b_i-),H_w(b_i+) = H_w(a_{i+1}-)$.


Following the claims above, for any action at $a \in (u+, t)$, we have $H_w(a+) \geq H_w(a-)$.

By the claim about \textbf{P} action, we have: 
$$
H_w(b_i-) - H_w(b_i+) \leq D(t(b_i))
$$
where $t(b_i)$ denote the time duration of $b_i$.

Therefore, we obtain:
\begin{align}
H_w(u+) - H_w(t-) &= \sum_{i=0}^k \left(H_w(a_i+) - H_w(a_{i+1}-)\right) + \sum_{i=1}^k \left(H_w(a_i-) - H_w(a_i+)\right)\\
&= \sum_{i=0}^k \left(H_w(b_i-) - H_w(b_i+)\right) + \sum_{i=1}^k \left(H_w(a_i-) - H_w(a_i+)\right)\\
&\leq \sum_{i=0}^k D(t(b_i)) = D(t-u) \leq D(30s)
\end{align}
where $a_0$ denote $u$ and $a_{k+1}$ denote $t$.

Therefore:
$$H_w(t-) > H_w(u+) - D(30s) \geq H_d - D(30s) \geq H_f
$$

This implies that $H(t-)\geq H_w(t-) >H_f$.

% Therefore, the decrease in $\frac{V}{IM}$ from $u$ to $t_0$ can only be caused by the change of rate. Since $t_0 - u< 30s$, it's guarantee that $V>\bt IM$ at $t_0-$. 

% WLOG, assume that all the strong orders are filled first. By Corollary 6.2, we know that after all these orders are filled, the health ratio of user cannot drop below $H_d$. After that, by Proposition 8.1, we know that the health ratio of users must increases after each of the weak orders are filled. Therefore, the user's health is larger than $H_d$ after all orders in $O$ are filled.

\textbf{Case 2: If $H_w(t_1) < H_d$:} Consider the time $t_2 = t_1 - 30s$. 

If $H_w(t_2) > H_c$, then similar to previous section, we have $H_w(t_1) \geq H_w(t_2) - D(t_2-t_1) \geq H_c - D(t_2-t_1) \geq H_d$, contradiction.

Otherwise, if $H_w(t_2) < H_c$ then all of user's orders are canceled by bot 2 at some $u < t_1$. 

\textbf{Case 2.1: If there is at least one new strong action from $u$ to $t_1$}. Assume the last one is at $s\in(u,t_1)$. Similar to previous section, we have $H_w(s) \geq H_c\RA H_w(t_1) \geq H_c - D(t_1-s) \geq H_d$, contradiction.

\textbf{Case 2.2: If there is no new strong action from $u$ to $t_1$}. Recall the ``worst order set" $O$. Since all user's orders are canceled at $u$, we have $O(u) = \emptyset$. From the claim above, \textbf{Weak} actions don't change this set. \textbf{Fill Weak} also don't change this set. Therefore, $O(t_1) = O(u) = \emptyset$, and hence $H(t_1) = H_w(t_1)$, which is a contradiction since $H(t_1) \geq H_d > H_w(t_1)$ by assumption.

We conclude that proof here.

\end{proof}

\newpage


\section{Settlement}

Note that when $t$ is updated (every 8 hours), the user's health also changes. We want to make sure that these jumps are controllable. Through out this section, we denote the rate for settlement $r_s$ and the time where the settlement happens $t_s$. 
\subsection{Before $t_T$}

Recall user's health 
$$
H = \frac{c + s\times r_m\times t}{b\times s\times r_m^+\times t^+}
$$
where $t^+$ denote $\max(t, t_T)$. 

We consider the case $t^+ = t$, i.e. $t > t_T$. 

User's health can only jump from $H_1$ down to $H_2$ if and only if $s(r_s-r_m) < 0$ (user experienced a loss) and 
$$
|r_s - r_m| > r_m^+\times b\times (H_1 + (H_1-H_2)\times\frac{t-\Delta t}{\Delta t})
$$
where $\frac{t-\Delta t}{\Delta t}$ is the number of settlement left. 

\subsection{After $t_T$}
We consider the case $t_T <t$. User's health can only jump from $H_1$ down to $H_2$ if and only if $s(r_s-r) < 0$ (user experienced a loss) and 
$$
|r_s - r_m| > r_m^+\times b\times (H_1-H_2)\times\frac{t_T}{\Delta t}
$$
which is slightly more likely to happen compared to the previous case. 

\subsection{Discussion}

Let's say we set $H_1 =1$, $H_2 = H_c\approx 0.8$, then the RHS is not smaller than
$$
r_m^+\times b\times (1 + 0.2\times\frac{t_T}{\Delta t})
$$

If we further set $b=0.25,t_T = 10 \Delta t$, then it becomes $75\%\times r_m^+$. 

What if we set $H_1 =1$, $H_2 = H_d\approx 0.6$, then it becomes:
$$
r_m^+\times b\times (1 + 0.4\times\frac{t_T}{\Delta t}) \approx 125\%\times r_m^+
$$

And if we set $H_1 =1$, $H_2 = H_f\approx 0.4$, then it becomes:
$$
r_m^+\times b\times (1 + 0.6\times\frac{t_T}{\Delta t}) \approx 150\%\times r_m^+
$$

There is actually a decent chance it can happen. For example: $I = 10\%$ and $r_m = 10\%$. It just needs to settles a negative rate of $-5\%$.



\subsection{Controlling}
Our advantage is that we can foresee the next settlement. Therefore, we can set up the risk mitigation corresponding to the upcoming settlement. 

\begin{assumption}
    We can see the next settlement rate 60s before it's settled.
\end{assumption}

\begin{prop}
    For any user:
    $$
    H_{s+} \geq H_{s-} - \frac{\Delta t |r_0 - r_m|}{b \cdot r_m^+ \cdot (t_s-\Delta t)^+}
    $$
    where $H_{s+}, H_{s-}$ are the user's health right before and right after the settlement.

    Specifically, the inequality happens when we observed an upcoming settlement with a loss for this user, i.e. $s(r_0-r_m) < 0$, and $t_s \leq T$.

    We call this jump "Settlement Health Jump", denoted by $D_s$:
    $$
    D_s(r_m, r Min_0) := \frac{\Delta t |r_0 - r_m|}{b \cdot r_m^+ \cdot (t_s-\Delta t)^+}
    $$
\end{prop}
\begin{proof}
When $t_s > T$:
$$
H_{s+} = H_{s-} + \frac{\Delta t (sign(s) \cdot (r_0-r_m) + H_{s-} \cdot b \cdot r_m^+)}{b \cdot r_m^+ \cdot (t_s - \Delta t)} > H_{s-} + \frac{\Delta t \cdot sign(s) \cdot (r_0-r_m)}{b \cdot r_m^+ \cdot (t_s - \Delta t)}
$$

When $t_s \leq T$:
$$
H_{s+}  = H_{s-} + \frac{sign(s) \cdot \Delta t (r_0 - r_m)}{b \cdot r_m^+ \cdot T}
$$

The argument follows immediately.
\end{proof}

\begin{remark}
    One may ask how big is $D_s$. An intuitive estimation would be when $t_s \leq T$, setting $T = 20$ (7 days) and $b = 0.25$, we have $D_s=\frac{|r_0 -r_m|}{5r_m^+}$ (e.g. 0.6 when $r_0 = 3r_m > I$). This value scales down linearly w.r.t. $t_s$ as it's further from maturity. 
\end{remark}

Recall that the worst health $H_w$ of user is less volatile than his Health. We have the following Corollary.

\begin{coro} For any user: 
$$
H_w(s_+) \geq H_w(s_-) - D_s
$$
\end{coro}

Now we're ready to take the settlement into account the settlement in Proposition \ref{safe-core-1}.

\begin{assumption}[Buffered bot when near to settlement]\label{bots-settle}
    Suppose we observed the settlement with rate $r_0$. During the period $(t_s-60s, t_s)$, Bot 1 and Bot 2 in Assumption \ref{bots} works with a buffer in their threshold:
    
    (i) Bot 1 is allowed to deleverage user when $H < H_d + \sup_{r(t_s)|r_m}D_s(r_0, r(t_s))$ \textcolor{red}{and his position size $s$ satisfies $s(r_s-r_m) < 0$}.
    
    (ii) Bot 2 is allowed to cancel user's open order when $H_w < H_c + \sup_{r(t_s)|r_m}D_s(r_0, r(t_s))$. 
\end{assumption}

\begin{prop}\label{safe-core-2}
    Consider time $t$. User's health always satisfies $H \geq H_f$. 
\end{prop}
\begin{proof}
    Apply the arguments in Proposition \ref{safe-core-1}, noticing that $H_w$ and $H$ cannot jump down more than $D_s$ by the settlement.
\end{proof}








% \begin{prop}
% At time $t$, user's collateral can tolerate the average rate up to $\frac{t_{th}\times k_{MM}\times\max\{{I,|r_m|}\}+t\times r_m}{t}$ given that he is not liquidatable.
% \end{prop}
% \begin{proof}
%     At time $t < t_{th}$, user needs to pay $t\times N_p$ more payments. Suppose that user $i$ is not liquidatable yet:
%     \begin{align*}
%     c &\geq  -s\times MM_m\times t_{th}-s\times t\times r_m\text{ for } s<0\\
%     c &\geq  s\times MM_m\times t_{th}-s\times t\times r_m\text{ for } s>0\\
%     \end{align*}

%     Suppose that the average rate of the payments until maturity is $r_a$, then the collateral covers these payments if and only if:
%     $$
%     c + s\times t\times N_p\times r_a \geq 0
%     $$

%     For $s < 0$, this is equivalent to:
%     \begin{align*}
%         s\times MM_m\times t_{th}+s\times t\times r_m &\leq s\times t\times r_a\\
%         \iff t_{th}\times MM_m+t\times r_m &\geq t\times r_a\\
%         \iff \frac{t_{th}\times MM_m+t\times r_m}{t} &\geq r_a
%     \end{align*}
%     which is true.
% \end{proof}
% \begin{remark} 
%     Let the market duration is 3 months, meaning that we have to scale up all the rate by 4 to get the annualized rate. 
    
%     Let $t_{th} = \frac{30}{3\times 365}$ which is 30 payments, or 10 days until maturity. Let $I = 2.5\%$ which corresponds with $I=10\%$ when the maturity is 3 months. Let $k_{MM} = 0.25$. 

%     At $t=5$ payments left, the collateral can cover about $16\%$ in average difference between $r_a$ and $r_m$.

%     At $t=2$ payments left, the collateral can cover about $37\%$ in average difference between $r_a$ and $r_m$.

%     At $t=1$ payments left, the collateral can cover about $75\%$ in average difference between $r_a$ and $r_m$.
% \end{remark}

\section{Appendix}

\subsection{Limit Order Range Derivation}\label{limit-range}
In this appendix, we derive the function 
$$f^u(r_m) := \min_x\{\max(g^{-1}(f_1(x;r_m)), f_2(x;r_m))\}$$ 
and 
$$f^l(r_m) := \max_x\{\min(h^{-1}(f_3(x;r_m)), f_4(x;r_m))\}$$ 
in Proposition \ref{core-limit-range}. We will show that:

\begin{equation*}
f^u(r; b,a) =
\begin{cases}
  2r\times\frac{1-b}{2-b-a} & r > I \\
  \frac{2(r-Ib)}{2-a-b} & I > r > I(1-\frac{a-b}{2}) \\
  r + \frac{(a-b)I}{2} & 0 < r <I(1-\frac{a-b}{2})
\end{cases}
\end{equation*}

\begin{equation*}
f^l(r; b,a) =
\begin{cases}
  r\times\frac{2+2b}{2+a+b} & r > I\frac{2+b+a}{2+2b} \\
  r(1+b) - \frac{(a+b)I}{2} & I < r < I\frac{2+b+a}{2+2b} \\
  r - \frac{(a-b)I}{2} & 0 < r <I
\end{cases}
\end{equation*}

For $f^u(r)$, recall the related functions:
\begin{align*}
    f_1(x;a,b,r) &= r + r^+(a x - b|1+x|)\\
    f_2(x;a,b,r) &= r + r^+(a|x| - b|1+x|)\\
    g(x) &= x - a\times x^+
\end{align*}

One can find that for any $r,a>0$ and $b<a$, $f_1(x)$ is a non-decreasing function, $f_2(x)$ is non-increasing on $(-\infty,0]$ and non-decreasing on $[0,\infty)$. Since $g(x)$ is a non-decreasing function for $(a<1)$, $g(f_2(x))$ is also non-increasing on $(-\infty,0]$ and non-decreasing on $[0,\infty)$.

Based on this structure of two functions $f_1$ and $g(f_2)$, we know that $\min_x\max(f_1(x),g(f_2(x)))$ (as well as $\min_x\{\max(g^{-1}(f_1(x;r_m)), f_2(x;r_m))\}$, since $g$ is non-decreasing) is obtained at one of intersections of $f_1(x)$ and $g(f_2(x))$. Moreover, since $f_1(x)$ is non-decreasing, the minimizer is always the negative intersection (if such intersection point exists). Note that there is at most 1 negative intersection point between $f_1(x)$ and $g(f_2(x))$.

Solving for a negative intersection between these two graphs, we obtain the minimizer $x^*$. Substitute this minimizer into $f_2$, we get the exact formula for $f^u(r) = f_2(x^*;r)$ as above.

\[
f^{u}(r) := \min_{x} \left\{ \max \left( g^{-1}(f_1(x; r)), \; f_2(x; r) \right) \right\}
\]\[
f_1(x; a, b, r) = r + r^{+}(ax - b|1 + x|)
\]
\[
f_2(x; a, b, r) = r + r^{+}(a|x| - b|1 + x|)
\]
\[
g(x) = x - a \times x^{+}
\]
One can find that for any $r, a > 0$ and $b < a$, $f_{1}(x)$ is a non-decreasing function, 
$f_{2}(x)$ is non-increasing on $(-\infty,0]$ and non-decreasing on $[0,\infty)$. 
Since $g(x)$ is a non-decreasing function for $(a < 1)$, $g(f_{2}(x))$ is also 
non-increasing on $(-\infty,0]$ and non-decreasing on $[0,\infty)$.

Based on this structure of two functions $f_{1}$ and $g(f_{2})$, we know that
\[
\min_{x} \max \big(f_{1}(x), g(f_{2}(x))\big)
\]
as well as 
\[
\min_{x}\left\{ \max \big(g^{-1}(f_{1}(x;r_m)), f_{2}(x;r_m)\big) \right\},
\]
since $g$ is non-decreasing) is obtained at one of intersections of $f_{1}(x)$ and $g(f_{2}(x))$. 
Moreover, since $f_{1}(x)$ is non-decreasing, the minimizer is always the negative intersection 
(if such intersection point exists). Note that there is at most 1 negative intersection point between 
$f_{1}(x)$ and $g(f_{2}(x))$.

Solving for a negative intersection between these two graphs, we obtain the minimizer $x^{*}$. 
Substitute this minimizer into $f_{2}$, we get the exact formula for
\[
f^{u}(r) = f_{2}(x^{*};r)
\]
as above. The same argument applies for $f^l(r)$. We derive it rigorously in the following subsections.


\subsection{ Upper bound}

$\textbf{Case 1}$ : $r>I \implies r^+ = r \implies f_1(x;r) = r +r(ax-b|1+x|)$ and $f_2(x;r) = r+ r(a|x| - b|1+x|) $\\ 
Solving $f_1(x;r) = g(f_2(x;r)) $ we have 
\[ r+ rax - rb|1+x| = r+ ra|x| - rb|1+x| - a \max\{I,f_2(x)\} \] 
\[\iff  rx = r|x| -  \max \{ I, f_2(x) \} \]
As we analyzed above, we go into case where $x < 0 $ first. If there exists a solution then it must be global minimum point.. 


\textbf{Case 1.1 }: $x<0$
\[2rx = - \max\{I,f_2(x) \}  \]

\textbf{ Case 1.1.1} : $x<-1$
\[ f_2(x) = r + rb + r(b-a)x \]
we have  $I/r -1 -b <a-b $ since $(0<I<r)$, hence $\frac{I/r -1 -b}{b-a}> -1 $ $\implies x<-1<\frac{I/r -1 -b}{b-a}  $\\ 
$\iff x(b-a)> I/r -1-b \iff r(1+b+(b-a)x) >I  \iff f_2(x)>I  $$\implies -\max\{I,f_2(x) \} = -f_2(x)$\\
Solving equation, we have $x^* = - \frac{1+b}{b+2-a} > -1 \implies$ this can't be solution 

\textbf{Case 1.1.2} : $\frac{I/r-1-b}{b-a} >x\ge -1$
we still have $\max\{ I,f_2(x) \} = f_2(x)= r-rb -r(a+b)x  \\ \implies x^* = \frac{1-b}{2-a-b} \implies f_2(x^*;r) = r - rb -r(a+b) \frac{1-b}{2-a-b} = 2r \frac{1-b}{2-a-b}$

\textbf{Case 1.1.3 } : $x \ge \frac{I/r-1-b}{b-a} > -1  \implies \max\{ I,f_2(x)\} = I  $ and $f_2(x) = r-rb-r(a+b)x\\  \implies x^* = \frac{-I}{2r} $. However, we can check that $\frac{-I}{2r} < \frac{I/r -1-b}{b-a} \iff I( \frac{2}{a-b}-1) \ge 2r \frac{1-b}{a-b} \implies I( \frac{2}{a-b}-1) \ge 2I \frac{1-b}{a-b} \iff a>b$, so this confirms there only exists an unique solution and this value can't be solution. \\ 
 $\implies $In this case, $f^u(r;b,a) = 2r \frac{1-b}{2-a-b}(*)$ 

$\textbf{Case 2}$ : $I>r> I( 1- \frac{a-b}{2}) \implies r^+ = I \implies f_1(x;r) = r +I (ax - b|1+x|) $ and $f_2(x;r) = r+I(a|x| - b|1+x|)$ 
Solving $f_1(x;r) = g(f_2(x;r)) $ we have 
\[ r+ I(ax - b|1+x|)  = r+ I(a|x| - b|1+x|) - a \max \{ I, f_2(x) \}\]
\[\iff  I x =I|x| -\max \{ I, f_2(x) \}  \]
As we analyzed above, we go into case where $x < 0 $ first. If there exists a solution then it must be global minimum point 

$\textbf{Case 2.1}:x<0$. This equation becomes : 
\[ 2Ix = - \max\{ I,f_2(x)\} \]
Assume for $x \in L : f_2(x) \le I \implies x^* = \frac{-1}{2}$, \\we check again $f_2(-\frac{1}{2}) = r - I \frac{b-a}{2 } > I - I \frac{a-b}{2} - I\frac{b-a}{2}  = I $ \\
This clearly violate assumption that $f_2(x) \le I$
So we want to find $x^* \notin L : f_2(x) > I  $\\
$\textbf{Case 2.1.1} : x<-1 :$
\[ f_2(x) = r + Ib+I(b-a)x  \implies x^* = \frac{r+Ib}{I(-2+a-b)} +1 -1  = \frac{r -I(2+b)}{I(-2+a-b)} -1 > -1  \]
this shows contradiction, hence there is no solution in this case

$\textbf{Case 2.1.2}: 0> x\ge -1$:
\[f_2(x) = r - Ib - I(a+b)x \implies x^* = \frac{r-Ib}{I(2-a-b)}  \]\[ \implies f_2(x^*) = r- Ib + I (b+a) \frac{r-Ib}{I(2-a-b)} =  \frac{2(r-Ib)}{2-a-b}\]

$\implies$ in this case, $f^u(r;b,a) = \frac{2(r-Ib)}{2-a-b}(**) $

\textbf{Case 3 }: $0<r<I(1-\frac{a-b}{2})  \implies r^+ = I \implies f_1(x;r) = r +I (ax - b|1+x|) $ and $f_2(x;r) = r+I(a|x| - b|1+x|)$ 
Solving $f_1(x;r) = g(f_2(x;r)) $ we have 
\[ r+ I(ax - b|1+x|)  = r+ I(a|x| - b|1+x|) - a \max \{ I, f_2(x) \}\]
\[\iff  I x =I|x| -\max \{ I, f_2(x) \}  \]

\textbf{Case 3.1} : $x \in L : f_2(x) \le I \implies  x^* = \frac{-1}{2} \implies f_2(x) = r + \frac{I(a-b)}{2}(***)$

\textbf{Case 3.2 } : $ x\notin L : f_2(x) > I$. Similarly to case $2.1.1$ and $2.1.2$, we also found that $f_2(x^*) =  \frac{2(r-Ib)}{2-a-b} < 2I( 1-b-\frac{a-b}{2}) / (2-a-b) = I$, this shows contradiction to our assumption, hence, in this case there is no solution

\textbf{Conclusion}: Combine $(*), (**),(***)$, we conclude that 
\[
f^u(r; b, a) =
\begin{cases}
  2r \times \dfrac{1-b}{2-b-a}, & r > I \\[1.2em]
  \dfrac{2(r - Ib)}{2-a-b}, & I > r > I\left(1 - \dfrac{a-b}{2}\right) \\[1.2em]
  r + \dfrac{(a-b)I}{2}, & 0 < r < I\left(1 - \dfrac{a-b}{2}\right)
\end{cases}
\]


\subsection{Lower bound}

Recall 
\[
f^{l}(r) := \max_{x} \left\{ \min \left( h^{-1}(f_3(x; r)), \; f_4(x; r) \right) \right\}
\]\[
f_3(x; a, b, r) = r - r^{+}(ax - b|1 + x|)
\]
\[
f_4(x; a, b, r) = r - r^{+}(a|x| - b|1 + x|)
\]
\[
h(x) = x + a \times x^{+}
\]
For any $r,a>0$ and $b<a$, $f_3(x)$ is non-increasing function, $f_4(x)$ is non-decreasing on $(-\infty ,0]$ and non-increasing on $[0,\infty)$. Since $h(x)$ is a non-decreasing function for $a>0$, $h(f_4(x))$ is also non-decreasing on $(-\infty ,0]$ and non-increasing on $[0,\infty)$. \\
Based on this structure of 2 function $f_3, h(f_4)$ we know that \[\max_x \min (f_3(x), h(f_4(x))  \]
as well as \[ \max_x \{ \min (h^{-1}(f_3(x;r), f_4(x;r))  \} \]
is ontained at one of intersection of $f_3(x)$ and $h(f_4(x))$. Moreover, since $f_3(x)$ is non-increasing, the maximizer is always the negative intersection ( if such intersection point exists). Note that there is at most 1 positive intersction point between $f_3(x) $ and $h(f_4(x))$. \\ 
Solving for a negative intersection, we obtain the maximizer $x^*$. Substitute this maximizer into $f_4$, we get the exact formula for 
\[ f^l(r) = f_4(x^*;r) \] as above 

\textbf{Case 1}: $0<r < I \implies r^+ = I \implies f_3(x) = r - I(ax-b|1+x|)$ and $f_4(x) = r - I(a|x| -b|1+x|)$.

We want to solve 
\[ r - I(ax-b|1+x|) = r - I(a|x| -b|1+x|) + a \times \max\{f_4(x),I \}   \]
\[\iff I|x| = Ix + \max\{f_4(x),I \} \]
Obviously we investigate $x<0: \implies  -2Ix = \max \{ f_4(x),I\}  $
From this we found $x^* = \frac{-1}{2} \implies f_4(x) = r -  I \frac{a-b}{2}$ satisfy this equation. And as analyze above, there is at most 1 negative intersection point, hence, in this case:
\[ f^l(r) = r - \frac{I(a-b)}{2} (*)  \]

\textbf{Case 2} : $I<r<I \frac{2+b+a}{2+2b} \implies r^+ = r \\ \implies f_3(x) = r - r(ax - b|1+x|)$ and $f_4(x) = r - r(a|x| - b|1+x|)$. We want to solve 
\[  r - r(ax - b|1+x|) = r - r(a|x| - b|1+x|) + a \max \{I,f_4(x) \}  \]
\[\iff r|x| = rx + \max \{I,f_4(x) \} \]
\textbf{Case 2.1} : $x<-1$: 
\[ f_4(x)/ r = 1 + ax - b(1+x) = (1-b) + (a-b)x \le (1-b)-(a-b) = 1-a \]
\[\implies f_4(x) \le r(1-a) \le \frac{2+b+a}{2+2b} (1-a)I  \le  I \frac{2+b+a - 2a}{2+2b} < I \]
So we have 
\[ -2rx = I \iff x = \frac{-I}{2r}  > -1     \]
which shows a contradiction to our assumption,

\textbf{Case 2.2} : $0> x \ge -1 \implies f_4(x)/r = 1+ax +b(1+x) = 1+b+ x(a+b) $

\textbf{Case 2.2.1} Again, assume for $x \in L : $$  f_4(x) > I $. This equation becomes \[ -2x = 1 + ax +b(1+x) \iff x = - \frac{1+b}{a+b+2}  \] \[ \implies f_4(x) =r(  1+b - (a+b) \frac{1+b}{a+b+2}) = \frac{2r(1+b)}{2+a+b}  < I  \]
since $r< \frac{2+a+b}{r(2+2b)}$. This violates our assumption. 

\textbf{Case 2.2.2 } Assume for $x \in L': f_4(x) \le I \implies x^* = \frac{-I}{2r} $\[ \implies f_4(x^*) = r(1+b) - I \frac{a+b}{2} \]
We can verify that $f_4(x^*) \le I$\\
So in this case \[ f^l(r) = r(1+b) - I \frac{a+b}{2} (**)\]

\textbf{Case 3:} $r> I \frac{2+b+a}{2+2b}$ 

Using the same analysis step as for the previous case, we only need to consider case : $ 0 >x \ge -1 $. Assume for $x \in L': f_4(x) \le I \implies x^* = \frac{-I}{2r} $\[ \implies f_4(x^*) = r(1+b) - I \frac{a+b}{2} > \frac{2+b+a}{2+2b} (1+b) I - I \frac{a+b}{2} = I\]
Hence, it violates our assumption.\\
Again, assume for $x \in L : $$  f_4(x) > I $. This equation becomes the following equation. the following equation. \[ -2x = 1 + ax +b(1+x) \iff x^* = - \frac{1+b}{a+b+2}  \] \[ \implies f_4(x^*) =r(  1+b - (a+b) \frac{1+b}{a+b+2}) = \frac{2r(1+b)}{2+a+b}  > I  \]
So in this case \[ f^l(r) = \frac{2r(1+b)}{2+a+b} (***) \]

Combined with $(*),(**),(***)$, we derive an exact formula. 
\[
f^{l}(r; b, a) =
\begin{cases}
r \times \dfrac{2+2b}{2+a+b}, & r > I \dfrac{2+b+a}{2+2b}, \\[1.2em]
r(1+b) - \dfrac{(a+b)I}{2}, & I < r < I \dfrac{2+b+a}{2+2b}, \\[1.2em]
r - \dfrac{(a-b)I}{2}, & 0 < r < I
\end{cases}
\]


\subsection{General setting}

For the mark price $-r < 0$: 
\begin{align*}
    f^u(-r) &= -f^l(r)\\
    f^l(-r) &= -f^u(r)
\end{align*}

We also want a simpler version of these bounds, $f^l_s$ and $f^u_s$ to be implemented in the contract. These bounds need to satisfy: $f^l_s > f^l$ and $f^u_s < f^u$, so when the price range requirement satisfy $f^l_s(r_m) < r < f^u_s(r_m)$, it also satisfies $f^l(r_m) < r < f^u(r_m)$, and hence the condition in Proposition 6.2 follows.

We propose a simpler version as following:
\begin{equation*}
f_s^u(r) =
\begin{cases}
  2r\times\frac{1-b}{2-b-a} & r > I \\
  r + \frac{(a-b)I}{2} & 0 < r <I
\end{cases}
\end{equation*}

\begin{equation*}
f_s^l(r) =
\begin{cases}
  r\times\frac{2+2b}{2+a+b} & r > I \\
  r - \frac{(a-b)I}{2} & 0 < r <I
\end{cases}
\end{equation*}
    


\end{document}